\documentclass[reqno]{amsart} 
\usepackage{amssymb,latexsym,amsmath,amscd,graphicx,setspace,amsthm,verbatim,enumitem}
\usepackage[margin = 3 cm]{geometry}

\def\upint{\mathchoice%
    {\mkern13mu\overline{\vphantom{\intop}\mkern7mu}\mkern-20mu}%
    {\mkern7mu\overline{\vphantom{\intop}\mkern7mu}\mkern-14mu}%
    {\mkern7mu\overline{\vphantom{\intop}\mkern7mu}\mkern-14mu}%
    {\mkern7mu\overline{\vphantom{\intop}\mkern7mu}\mkern-14mu}%
  \int}
\def\lowint{\mkern3mu\underline{\vphantom{\intop}\mkern7mu}\mkern-10mu\int}


\theoremstyle{plain}
\newtheorem{theorem}{Theorem}[section]
\newtheorem{proposition}{Proposition}
\newtheorem{corollary}{Corollary}
\newtheorem{lemma}{Lemma}
\newtheorem{conjecture}{Conjecture}
\newtheorem{question}{Question}
\newtheorem{problem}{Problem}
      
\theoremstyle{definition}
\newtheorem{definition}{Definition}

\newenvironment{solution1}{\paragraph{\emph{Solution $1$}.}}{\hfill$\square$}
\newenvironment{solution2}{\paragraph{\emph{Solution $2$}.}}{\hfill$\square$}
\newenvironment{solution3}{\paragraph{\emph{Solution $3$}.}}{\hfill$\square$}

\begin{document} 

\title[Homework 2]{Homework 2}

\date{\today} 
\maketitle 
 

\begin{problem}[{\bfseries optional} for bonus points]
Show that if $f:[a,b] \rightarrow \mathbb{R}$ is a bounded function that is integrable in the sense of Darboux, then it is integrable in the sense of Riemann.
\end{problem}
\begin{solution1}

\end{solution1}

Since we proved the converse in class, we now know that Darboux integrability is equivalent to Riemann integrability for real-valued functions bounded on a closed interval $[a,b]$.


\begin{problem}[1B-3]
Suppose $f,g:[a,b] \rightarrow \mathbb{R}$ are bounded functions.  Prove that
$$\lowint_{a}^{b}f + \lowint_{a}^{b}g \le \lowint_{a}^{b}(f+g) $$
and
$$\upint_{a}^{b}(f+g) \le \upint_{a}^{b}f + \upint_{a}^{b}g  $$
\end{problem}
\begin{solution1}

\end{solution1}

\begin{problem}
Consider the Dirichlet function $f:[0,1] \rightarrow \mathbb{R}$ defined via
$$x \mapsto f(x) = \begin{cases}1, &\text{ if } x \in \mathbb{Q}\cap [0,1];\\ 0, &\text{ if } x \notin \mathbb{Q}\cap[0,1].  \end{cases} $$
Show that $f$ is not Riemann integrable.
\end{problem}
\begin{solution1}

\end{solution1}


\begin{problem}[1B-4]
Give an example of bounded functions $f,g:[0,1] \rightarrow \mathbb{R}$ such that
$$\lowint_{0}^{1}f + \lowint_{0}^{1}g < \lowint_{0}^{1}(f+g) $$
and
$$\upint_{0}^{1}(f+g) < \upint_{0}^{1}f + \upint_{0}^{1}g  $$
\end{problem}
\begin{solution1}

\end{solution1}


\begin{problem}[1B-5]
Give an example of a sequence of continuous real-valued functions $f_{1},f_{2},\ldots$ on $[0,1]$ and a continuous real-valued function $f$ on $[0,1]$ such that
$$f(x) = \lim_{k \to \infty} f_{k}(x) $$
for each $x \in [0,1]$ but
$$\int_{0}^{1}f \neq \lim_{k \to \infty} \int_{0}^{1} f_{k}. $$
\end{problem}
\begin{solution1}

\end{solution1}

\begin{problem}[1B-1]
Consider the Thomae function $f:[0,1] \rightarrow \mathbb{R}$ defined via
$$x \mapsto f(x) = \begin{cases}\frac{1}{n}, &\text{ if } x = \frac{m}{n} \in \mathbb{Q}\cap [0,1], \text{ where } \frac{m}{n} \text{ is in reduced form and } n > 0;\\ 0, &\text{ if } x \notin \mathbb{Q}\cap[0,1].  \end{cases} $$
Show that $f$ is Riemann integrable and compute $\int_{0}^{1} f$.
\end{problem}
\begin{solution1}

\end{solution1}



In class, we defined the Lebesgue outer measure of a subset of $\mathbb{R}^{n}$ in terms of {\bfseries closed} $n$-dimensional boxes.  But one could use {\bfseries open} $n$-dimensional boxes instead as presented in the book.  An {\bfseries open} $n$-dimensional box is a subset of $\mathbb{R}^{n}$ of the form
$$I = (a_{1},b_{1}) \times (a_{2},b_{2}) \times \ldots \times (a_{n},b_{n}), $$
and one defines its $n$-dimensional volume to be the same as the corresponding {\bfseries closed} box when the endpoints of each interval are included.

\begin{problem}
For the purpose of this problem let $n=1$ and for $A \subseteq \mathbb{R}$, define
$$\lambda^{0}(A) = {\rm inf}\left\{\sum_{k=1}^{\infty}\ell(I_{k}) : I_{1},I_{2},\ldots \text{ is a sequence of open intervals satisfying } A \subseteq \bigcup_{k=1}^{\infty} I_{k}  \right\}. $$
Show that $\lambda^{0}(A) = \lambda^{*}(A)$ for all subsets $A$ of $\mathbb{R}$.  (A similar statement is true for $n > 1$, but I leave that to you to think about on your own.)
\end{problem}
\begin{solution1}

\end{solution1}




\end{document}









